\begin{theorem}[Unramified case]\label{mod:cases-unramified}
Let $K/F$ be a prime cyclic extension of nonarchimedean local fields with
$e(K/F)=1$, let $\chi_F$ and $\psi_F$ carry proofs of their exact
multiplicative and additive conductors, and choose a uniformizer $\pi$ of
$F$.  Write
\[
 \chi_K=\chi_F\circ N_{K/F},\qquad
 \psi_K=\psi_F\circ\operatorname{Tr}_{K/F},\qquad
 \gamma_F=\pi^{m_F(\chi_F)+n_F(\psi_F)},
\]
and let $\gamma_K$ be the image of $\gamma_F$ in $K^\times$.

For every $\mu\in S(K/F)$, Lean packages $\mu$ with exact conductor zero
and $\mu\chi_F$ with the unchanged exact conductor of $\chi_F$.  With the
same uniformizer denominator in every factor, it proves the complete
computational First Main identity
\[
 \Delta_K^{\mathrm{fin}}(\chi_K,\psi_K;\gamma_K)
 \prod_{\mu\in S(K/F)}
   \Delta_F^{\mathrm{fin}}(\mu,\psi_F;\gamma_\mu)
 =
 \prod_{\mu\in S(K/F)}
   \Delta_F^{\mathrm{fin}}(\mu\chi_F,\psi_F;\gamma_{\mu\chi_F}).
\]
Here $\Delta^{\mathrm{fin}}$ is Lean's \texttt{deltaFinite} with the
project's normalization, and the products include the trivial norm
character.

Lean also exports the dispatch-ready wrapper
\texttt{firstMain\_unramified\_dispatchReady}.  Given proposed local-constant
functions $\Delta_F,\Delta_K$ satisfying
\texttt{IsDeltaFiniteLocalConstant}, it constructs
\texttt{FirstMainComputationalData} from the same private unramified
conductor packages, chooses the same uniformizer denominators, and identifies
\texttt{unramifiedNormCharacterFinset} with the full norm-character
enumeration, including the trivial character.  It supplies the displayed
computational identity from \texttt{firstMain\_unramified} to
\texttt{firstMainIdentity\_of\_deltaFinite\_product}, yielding
\[
 \texttt{FirstMainIdentity}(F,K,\Delta_F,\Delta_K,
   \chi_F,\psi_F).
\]
Thus the wrapper adds only the dispatch interface; it neither repeats the
finite-field argument nor duplicates error-term cancellation.

The proof first enumerates all norm characters by
$\mu\mapsto\mu(\pi)$ and proves
\[
 \prod_{\mu\in S(K/F)}\mu(\pi)=(-1)^{[K:F]-1}.
\]
It then reduces the displayed identity to the exact signed power formula
\[
 \Delta_K^{\mathrm{fin}}(\chi_K,\psi_K;\gamma_K)
 =(-1)^{m_F(\chi_F)([K:F]-1)}
  \Delta_F^{\mathrm{fin}}(\chi_F,\psi_F;\gamma_F)^{[K:F]}.
\]
Conductor zero follows from the unramified-character formula.  At conductor
one, the residue multiplicative character is pulled back by the norm and
the residue additive character by the trace; Hasse--Davenport lifting is
applied with the existing inverse-character and leading-minus conventions.
For higher conductor, Lean splits $m=2d+\varepsilon$.  It first transports
the stationary numerator as an equality in the literal quotient
$\mathcal O_K/\mathfrak p_K^d$, and only then chooses mapped integral
representatives for Lamprecht's formula.  The even case has no residual
factor.  In the odd case the norm expansion retains the unordered second
coefficient, reduces it to the finite-field pair correction, identifies the
upstairs Hasse function with the trace lift carrying the exact negative
correction, and applies the Hasse-function lift.  Thus all three conductor
ranges, every exponent, and every sign are proved without a
Hasse--Davenport multiplication argument.
\LeanFile{LanglandsFirstMainLemma/Cases/Unramified.lean}
\lean{LanglandsFirstMainLemma.firstMain_unramified}
\lean{LanglandsFirstMainLemma.firstMain_unramified_dispatchReady}
\leanok
\SourceText{Proposition prop:unramified-new.}
\uses{mod:ramification-unramifiedcompatibility,mod:delta-residueformula,mod:finitefield-hassedavenportlift,mod:finitefield-hassefunctionlift,mod:lamprecht-formula,mod:parameters-highunramifiedstable,mod:ramification-normcharacters,mod:delta-firstmainstatement,mod:delta-errorterm}
\FormalizationContract The principal Lean theorem is the full
\texttt{deltaFinite} product identity for supplied exact conductor data and
a chosen uniformizer.  The dispatch-ready wrapper accepts computationally
compatible parameters of type \texttt{LocalConstantFunction}, packages the
same exact character data and denominators, normalizes to the full
\texttt{NormCharacter} product, and applies the generic product bridge to
prove \texttt{FirstMainIdentity}.  Every auxiliary theorem in the Lean file
remains private.
\end{theorem}
